%! Symmetric Positive Definite Matrices — Unit 6, Lesson 6.3 — Guided Notes
\documentclass[10pt]{article}
\usepackage{linalg-article}
\usepackage{linalg-boxes}
\newcommand{\TallMath}[1]{$\displaystyle #1\rule[-1.4em]{0pt}{3.2em}$}
\newcommand{\xx}{\mathbf{x}}
\newcommand{\vv}{\mathbf{v}}
\newcommand{\qq}{\mathbf{q}}
\newcommand{\zero}{\mathbf{0}}
\newcommand{\T}{\mathsf{T}}

\begin{document}

\pageheader{Unit 6, Lesson 6.3}{Guided Notes}
\namedateperiod

\vspace{0.12in}

\begin{objectivebox}
\textbf{By the end of this lesson, I will be able to\dots}
\begin{itemize}\setlength{\itemsep}{2pt}
  \item recognize a \emph{symmetric} matrix ($S=S^{\T}$) and explain why its eigenvectors are perpendicular;
  \item diagonalize a symmetric matrix as $S=Q\Lambda Q^{\T}$ with an orthonormal $Q$ (so $Q^{-1}=Q^{\T}$);
  \item test whether a symmetric matrix is \emph{positive definite}, and connect that to positive energy $\xx^{\T}S\xx>0$.
\end{itemize}
\end{objectivebox}

\vspace{0.06in}

\begin{vocabbox}
\small
Fill in each term as we build it together.
\vspace{2pt}
\termblanklong{Symmetric matrix ($S=S^{\T}$)}
\termblanklong{Orthonormal matrix $Q$ (perpendicular unit columns, $Q^{\T}Q=I$)}
\termblanklong{Positive definite (symmetric, all $\lambda>0$)}
\termblanklong{Energy $\xx^{\T}S\xx$}
\end{vocabbox}

\begin{hookbox}
In Lesson~6.2 a matrix diagonalized \emph{only} when it had enough independent eigenvectors --- and finding
$X^{-1}$ could be messy.
\begin{itemize}\setlength{\itemsep}{1pt}
  \item Is there a family of matrices that \emph{always} diagonalizes? \writeline
  \item Could the inverse $X^{-1}$ ever be \emph{free} to write down? \writeline
\end{itemize}
Today's answer to both: \textbf{symmetric} matrices, where $S=S^{\T}$.
\end{hookbox}

\begin{notesbox}{1. Symmetric Matrices Have Perpendicular Eigenvectors}
A matrix is \textbf{symmetric} when it equals its own transpose, $S=S^{\T}$ --- the numbers mirror across
the diagonal. Our unit matrix $S=\begin{bsmallmatrix}2&1\\1&2\end{bsmallmatrix}$ is symmetric (the two
off-diagonal $1$'s match). From Lesson~6.1 its eigen-data is
\[
  \lambda=3,\ \xx_1=\begin{bmatrix}1\\1\end{bmatrix};\qquad \lambda=1,\ \xx_2=\begin{bmatrix}1\\-1\end{bmatrix}.
\]
Check the angle between the eigenvectors with a dot product:
$\ \xx_1\cdot\xx_2=(1)(1)+(1)(-1)=\blank{1.0cm}$, so they are \blank{2.4cm}.

\vspace{2pt}
This is no accident. The \textbf{spectral theorem} says: \emph{every} symmetric matrix has all \textbf{real}
eigenvalues and \textbf{perpendicular} eigenvectors --- guaranteed, with no exceptions.

\begin{center}
\begin{tikzpicture}[scale=1.05, font=\scriptsize, >=stealth]
  % axes
  \draw[linegray, ->] (-2.2,0) -- (2.2,0) node[right]{$x$};
  \draw[linegray, ->] (0,-2.2) -- (0,2.2) node[above]{$y$};
  % eigen-lines (dashed)
  \draw[charcoal, dashed] (-1.9,-1.9) -- (1.9,1.9);
  \draw[charcoal, dashed] (-1.9,1.9) -- (1.9,-1.9);
  % right-angle mark at origin
  \draw[charcoal, thick] (0.34,0.34) -- (0.68,0) -- (0.34,-0.34);
  % eigenvectors
  \draw[burgundy, ultra thick, ->] (0,0) -- (1.6,1.6)
      node[above right, burgundy]{$\xx_1=\begin{bsmallmatrix}1\\1\end{bsmallmatrix}$ \ ($\lambda=3$)};
  \draw[royalblue, ultra thick, ->] (0,0) -- (1.6,-1.6)
      node[below right, royalblue]{$\xx_2=\begin{bsmallmatrix}1\\-1\end{bsmallmatrix}$ \ ($\lambda=1$)};
  \node[charcoal] at (0,-2.6) {\itshape symmetric $\Rightarrow$ eigenvectors are perpendicular};
\end{tikzpicture}
\end{center}
\end{notesbox}

\begin{notesbox}{2. The Cleanest Diagonalization: $S=Q\Lambda Q^{\T}$}
Perpendicular is good; \textbf{perpendicular and unit length} is even better. Normalize each eigenvector by
dividing by its length $\sqrt2$, and put the results in the columns of $Q$:
\[
  Q=\frac{1}{\sqrt2}\begin{bmatrix}1&1\\1&-1\end{bmatrix},\qquad
  \Lambda=\begin{bmatrix}3&0\\0&1\end{bmatrix}.
\]
Because the columns of $Q$ are perpendicular unit vectors, $Q^{\T}Q=I$ (Unit~4) --- so the inverse is just
the \textbf{transpose}: $Q^{-1}=\blank{1.2cm}$. The Lesson~6.2 formula $A=X\Lambda X^{-1}$ then becomes
\[
  S=Q\Lambda Q^{\T}\qquad(\text{no inverse to compute --- just flip }Q).
\]
Multiply it out to confirm (the $\frac1{\sqrt2}$'s combine to $\frac12$):
\[
  Q\Lambda Q^{\T}=\frac12\begin{bmatrix}3&1\\3&-1\end{bmatrix}\begin{bmatrix}1&1\\1&-1\end{bmatrix}
  =\frac12\begin{bmatrix}\rule{0pt}{1.1em}\blank{0.7cm}&\blank{0.7cm}\\[3pt]\blank{0.7cm}&\blank{0.7cm}\end{bmatrix}
  =\begin{bmatrix}2&1\\1&2\end{bmatrix}=S.\quad\checkmark
\]
\end{notesbox}

\begin{notesbox}{3. Positive Definite: All Eigenvalues Positive}
A symmetric matrix is \textbf{positive definite} when \emph{every} eigenvalue is positive. Our spine has
$\lambda=3$ and $\lambda=1$ --- both positive, so $S$ is positive \blank{1.8cm}.

\vspace{2pt}
For a $2\times2$ symmetric $S=\begin{bsmallmatrix}a&b\\b&c\end{bsmallmatrix}$ you don't even need the
eigenvalues --- there is a \textbf{quick test}:
\[
  S\text{ is positive definite}\iff a>0 \ \text{ and }\ ac-b^2>0.
\]
For the spine, $a=2>0$ and $ac-b^2=(2)(2)-1^2=\blank{1.0cm}>0$, so it passes. Notice $ac-b^2=\det S$, so a
positive-definite matrix always has $\det>0$ --- and therefore is \blank{2.2cm} (no zero eigenvalue).
\end{notesbox}

\begin{notesbox}{4. Energy $\xx^{\T}S\xx$: a Bowl That Curves Up}
There is one more way to see ``positive definite,'' and it is the most useful. The \textbf{energy} of a
vector $\xx=\begin{bsmallmatrix}x\\y\end{bsmallmatrix}$ is the number $\xx^{\T}S\xx$. For the spine,
\[
  \xx^{\T}S\xx = 2x^2+2xy+2y^2 = x^2+y^2+(x+y)^2.
\]
Written as a sum of squares, this is $>0$ for every $\xx\ne\zero$. That is exactly what positive definite
means: \textbf{$\xx^{\T}S\xx>0$ for all $\xx\ne\zero$}. Picture the surface $z=\xx^{\T}S\xx$ as a
\emph{bowl} sitting at the origin and curving up in every direction --- lowest point at $\zero$. A negative
eigenvalue would bend one direction \emph{down}, making a saddle instead of a bowl.

\vspace{3pt}
\textbf{The road ahead.} \textbf{6.4} puts eigenvalues to work one more time --- using
$A^k=X\Lambda^k X^{-1}$ to run a real changing system forward in time.
\end{notesbox}

\begin{practicebox}
Show each step.
\begin{enumerate}[leftmargin=1.6em, itemsep=6pt]
  \item $S=\begin{bsmallmatrix}3&1\\1&3\end{bsmallmatrix}$ has eigenvalues $\lambda=4,2$ with eigenvectors $\begin{bsmallmatrix}1\\1\end{bsmallmatrix},\begin{bsmallmatrix}1\\-1\end{bsmallmatrix}$. Confirm the eigenvectors are perpendicular, and use the quick test to show $S$ is positive definite. \writelines{2}
  \item Which of these symmetric matrices is positive definite? Use $a>0$ and $ac-b^2>0$. \quad $\begin{bsmallmatrix}2&1\\1&2\end{bsmallmatrix}$ \quad and \quad $\begin{bsmallmatrix}1&3\\3&1\end{bsmallmatrix}$. \writeline
  \item Why is the inverse of the orthonormal matrix $Q$ so easy to write down --- what makes $Q^{-1}=Q^{\T}$? \writeline
\end{enumerate}
\end{practicebox}

\end{document}
